% Abstract: a single paragraph, under 300 words. Instruction only.

[PROCESSING INSTRUCTION (Abstract, keep under 300 words, single
paragraph). Write one overview paragraph that moves through the full
pipeline in order: (1) bracketed instruction generation, (2) code
generation, (3) code execution, (4) chart or figure generation, and (5)
draft and then full paper production, all performed autonomously by
Claude Code Opus 4.7 (1M context) Max. State the thesis plainly, that
the LLM VVUQ process is the substantial, decision bearing part and must
be held higher than the generation it gates, and that this yields
faster, less expensive, and more rigorous physical AI oncology trials
than conventional verification. Close the paragraph with the headline
metrics, lightly woven (not a list): 51 of 51 tests passed across 8
modules; a 2.5x schedule acceleration (30 baseline to 12 automated
days); the VVUQ gate exercised across 6 cases that accepted 1 and
blocked 5 (1 of which escalated); a verification fraction of 1.0
required to ship; a maximum coefficient of variation of 0.0068 against a
0.10 bound; and the Stage 2 2030 PDAC analog (a 60 second 8 arm robotic
Whipple plus six 28 day adjuvant cycles, 168 regimen days).

SOURCE FILES to synthesize from (do not copy; abbreviate any path used
in prose):
- execution/README.md, the "Key quantitative results to cite" and
  "Execution results summary" tables, for every metric above;
- pipeline/README.md, vvuq/README.md, physical-ai/README.md, for the
  one line each on generation, assurance, and the Stage 2 reference.

HOW TO PROCESS: lead with motivation and thesis, then the five step
pipeline, then assurance, then the metrics; report numbers exactly as in
execution/README.md. Do not cite in the abstract. Keep it publication
grade and self contained.]
